\section{Development Methodology}
\label{sec:metodolohiia}

\subsection{Specification-Driven Development (SDD)}

The system is developed in accordance with the Specification Driven Development (SDD) methodology. The specification is the controlling source of system behavior --- no source code may define behavior independently of this document. \sked{K53}

Required process for every behavioral change:

\begin{enumerate}
  \item Update this document before making any code changes.
  \item Assign or update a requirement identifier (REQ-F-NNN or REQ-NF-NNN).
  \item Define an acceptance criterion (AC-NNN) for each requirement.
  \item Write an automated test before implementation (TDD, \secref{sec:metodolohiia}).
  \item Run the test and confirm it fails for the expected reason (Red).
  \item Implement the minimal code that passes the test (Green).
  \item Refactor only after tests pass (Refactor).
  \item Update log artifacts after completing the iteration.
\end{enumerate}

\subsection{Mandatory Test-Driven Development (TDD)}

TDD is mandatory for all implementation work. The Red~/ Green~/ Refactor cycle applies to every functional change: \sked{K53}

\begin{itemize}
  \item \textbf{Red} --- write a failing test first.
  \item \textbf{Green} --- implement the minimal code that passes the test.
  \item \textbf{Refactor} --- improve the code without changing external behavior.
\end{itemize}

No production code is accepted unless it is covered by at least one automated test. Real external systems (DBMS, third-party APIs) are not used in unit tests --- deterministic substitutes (stubs~/ mocks) are used instead.

\subsection{Traceability Rule}

Every implemented system behavior must have full traceability in the chain:

\begin{codelisting}
\caption{Traceability Chain}
\label{lst:trace}
\begin{skedcode}
REQ -> AC -> Test -> Implementation
\end{skedcode}
\end{codelisting}

Functionality without a requirement identifier is not part of the system. A requirement without an acceptance criterion is not ready for implementation.

\subsection{Test Log Artifacts}

After each test run, the development agent records a log file with the results. \sked{K54} The log file must contain:

\begin{itemize}
  \item total number of tests and the count of passed~/ failed;
  \item the status of each individual test (passed~/ failed~/ skipped);
  \item execution time of the test suite.
\end{itemize}

The log file is a machine-readable verification artifact --- based on it, the development agent and human reviewer determine whether the task has been completed according to the specification (the ``Is~verified?'' step in the SDLC process).
